\subsection{Maximal Submodules}

DEFINITION 2. --- Let A be a ring and M an A-module. A maximal submodule of M is a maximal element, for the inclusion, of the set of proper submodules of M.

A maximal submodule of \(A_{s}\) is simply a left maximal ideal of A.

Let N be a submodule of M. The submodules of M/N are of the form P/N, where P is a submodule of M containing N (I, §4, No.~6, p.~41, Theorem 4). Consequently, N is a maximal submodule of M if and only if the module M/N is simple.

PROPOSITION 3. --- Let M be a finitely generated A-module. Every proper submodule of M is contained in a maximal submodule.

Let N be a proper submodule of M. Denote by S the set of proper submodules of M containing N, ordered by inclusion; let us prove that S is inductive. Let F be a totally ordered subset of S. If F is empty, then N is an upper bound for F in S. In the opposite case, denote by Q the union of the elements of F. Then Q is a submodule of M. Let F be a finite generating subset of M. If Q were equal to M, then F would be contained in a submodule \(P \in F\) , which would imply P = M, contrary to the definition of F. We therefore have \(Q \in S\) , which proves that S is inductive. Proposition 3 then follows from Corollary 1 of Set Theory, III, §2, No.~4, p.~155.

COROLLARY 1. --- Let M be a nonzero finitely generated A-module. There exists a two-sided ideal a of A, annihilator of a simple A-module, such that aM is not equal to M.

Let N be a maximal submodule of M (Proposition 3), and let a be the annihilator of the simple A-module M/N; it is a two-sided ideal of A, and we have \(\mathfrak{a}(M/N)=0\) , whence \(aM\subset N\) and consequently \(aM\neq M\) .

COROLLARY 2. --- Let A be a commutative ring and B an A-algebra. Let M be a simple B-module that is a finitely generated A-module, and let m be the annihilator of the A-module M. Then m is a maximal ideal of A, and M is a finite-dimensional vector space over the field A/m.

Since the ring A is commutative, the annihilator of a simple A-module is a maximal ideal of A (VIII, p.~46, Proposition 1). By Corollary 1 applied to the A-module M, there exists a maximal ideal a of A such that \(aM \neq M\) . But aM is a submodule of the simple B-module M; we therefore have aM = 0 and consequently \(a \subset m\) . Since the ideal m is distinct from A, it is equal to a, so that m is maximal. The module M is a finitely generated module over the field A/m; the last assertion follows.

